\chapter{Moral Reasons: Impartiality, Deliberation and the Evaluation of Desire}
\label{sec:p11-rationality-12-2}

Michael Smith's 2021 chapter asks how an empirical psychology can make
room for moral reasons without leaving its philosophical assumptions
unspoken. The question has two parts: what distinguishes morality from
other practical standards, and what makes a consideration a reason for
action? Answers to the first do not automatically settle the second.
The source compares Rationalism and Antirationalism about morality with
Humeanism, Reasons Primitivism and Constitutivism about reasons. It
criticizes some combinations while explicitly leaving the merits of its
controversial premises open to assessment
\parencite[pp.~693--700]{Smith2021MoralReasons}.

This is a philosophical treatment, not an experiment establishing how
many people possess moral reasons or how often they act on them. Its
central distinction is between actual motivation and the motivation an
agent would have after deliberating well under specified conditions.
That counterfactual itself depends on an account of rational belief and
desire. A psychological theory cannot simply label observed helping,
cooperation or blame as evidence for one account without stating the
connection it assumes.

\section{The mark and the substance of morality}
\label{sec:p11-r122-morality}

\begin{hbdefinition}{def:p11-r122-impartiality}{Impartiality as the chapter's moral starting point}
Smith characterizes moral principles as practical principles requiring
impartial conduct, among other features. The moral standpoint considers
oneself as one among others entitled to the same consideration, rather
than assigning fundamental priority to oneself merely because one is
oneself. This characterization does not yet identify the scope of moral
consideration or the correct interpretation of impartiality.

Questions of scope concern who belongs to the relevant class: all
sentient beings, all rational agents, everyone to whom conduct must be
justifiable, or another specified group. Questions of substance concern
what equal consideration demands: equal weight for comparable interests
in promoting welfare, equal freedom to choose a life while respecting
others' freedom, or justification to each affected party. These are
different candidate interpretations, not one interchangeable decision
rule \parencite[pp.~693--694]{Smith2021MoralReasons}.
\end{hbdefinition}

The source therefore treats the impartiality characterization as compatible
with utilitarian, deontological and contractualist proposals. It invokes
Gauthier, Hare and Scanlon in introducing that broad moral standpoint,
not as proving that those theories have identical foundations. Its
exclusion is more specific: a fundamental principle giving a particular
person's welfare, freedom or recognition greater importance simply
because it is that person's own is not a moral principle on this
characterization. Smith accordingly presents thorough selfishness as a
paradigm of immorality.

That starting point is not a complete moral theory. It does not specify
how duties toward particular people are justified, how conflicting claims
are resolved, or what institutional rules ought to be adopted. Nor does
impartial consideration necessarily require identical treatment regardless
of differences in circumstances. The chapter's claim of broad agreement
on impartiality is an attributed conceptual starting point, not a measured
consensus or an independent argument settling every account of morality
\parencite[pp.~693--694]{Smith2021MoralReasons}.

The source contrasts this standpoint with the familiar Humean account
of rational belief and motivation. The contrast is not between caring
about oneself and never doing so: morality includes oneself among those
considered. It is between a standard whose content is impartial and a
practical-rationality standard that initially takes each agent's own
intrinsic desires as given, whatever their content.

\section{Humean rationality and the roles of belief and desire}
\label{sec:p11-r122-humeanism}

\begin{hbdefinition}{def:p11-r122-desires}{Intrinsic and instrumental desires in the source's Humean account}
An intrinsic desire concerns an outcome for its own sake. It can vary
in strength and need not be selfish: an agent can intrinsically desire
another person's well-being or impartial conduct. An instrumental belief
concerns which available option would satisfy an intrinsic desire; its
confidence can vary. An instrumental desire combines such intrinsic
desire and instrumental belief in an action-guiding motivation.

On the Humean account selected by Smith, epistemic principles assess
belief against available evidence. Practical principles assess motivation
given the intrinsic desires and instrumental beliefs the agent has.
They do not assess the intrinsic desires themselves as rational or
irrational. The source calls these principles self-referential because
the agent's own evidential and motivational starting points determine
the assessment, not because every agent must desire only personal gain
\parencite[pp.~694--695]{Smith2021MoralReasons}.

Degree of epistemic and practical rationality, on this account, concerns
having and exercising capacities to form beliefs and motivations in
conformity with these principles. It is not defined by the success of
one observed act alone.
\end{hbdefinition}

Smith reconstructs the Humean difference between the states as follows.
Both beliefs and desires have representational content, but a belief
represents something as being the case. An intrinsic desire represents
a possible way things could be, associated with positive affect or a
disposition to bring that condition about if one believed it available.
An intrinsic desire does not, merely by being a desire, assert its
content as true. Humeans infer from this contrast that intrinsic desires
are not sensitive to truth-supporting evidence in the way beliefs are,
and deny rational principles for evaluating their contents.

Instrumental desires differ because their belief component is assessable.
Evidence can show that the supposed means is unavailable or does not
produce the desired outcome. An instrumental motivation can therefore
be criticized while the intrinsic end remains uncriticized under this
account. The chapter's later alternatives target the transition from
``not truth-assessable'' to ``not rationally assessable,'' rather than
simply treating a desire as a false belief
\parencite[pp.~694--697]{Smith2021MoralReasons}.

The source associates probability-based belief revision and expected-
utility decision modeling with codifications of the Humean principles.
This is an interpretation of their use in psychology, not a theorem
that every probability model or every expected-utility representation
commits its user to Humeanism. The finite representation at
\cref{thm:p17-finite-vnm} establishes a numerical form from preference
axioms; it does not establish which desires are rationally permissible.
Likewise, the economic comparison at
\cref{sec:p19-rationality-9-4} distinguishes consistent preferences,
their numerical representation, cognitive processes and auxiliary
assumptions. Altruistic and morally loaded preferences are mathematically
possible without their moral authority being proved by the mathematics.

\section{Well-deliberated motivation: four distinctions}
\label{sec:p11-r122-deliberation}

\begin{hbdefinition}{def:p11-r122-deliberation}{The chapter's platitude and its idealizing conditions}
Smith takes as a common starting point that an agent has a reason to
do what the agent would be motivated to do if the agent deliberated
well. In the selected presentation, deliberating well involves three
conditions: beliefs conform to the principles of epistemic rationality;
motivations conform to the principles of practical rationality; and
both conform under circumstances in which relevant features of the
world are maximally and reliably revealed.

The first two conditions define evidence-relative, subjective reasons.
All three define fact-relative, objective reasons. The third is not
merely greater confidence: it excludes misleading disclosure and requires
the relevant facts to be reliably available. What those relevant facts
include depends on the account of rationality and reasons being used
\parencite[p.~695]{Smith2021MoralReasons}.
\end{hbdefinition}

Two further contrasts matter. First, an agent's actual motivation can
depart from what the agent would be motivated to do under those conditions.
Reason possession therefore does not require actual performance or even
the actual motivation predicted by the idealization. Second, having some
reason and having most reason are different. The source connects some
reason with any motivation that would survive good deliberation, whether
weak or strong, and most reason with the strongest such motivation.
It does not supply a quantitative calculus of strength, a universal
tie-breaking rule, or a measurement instrument for the counterfactual.

Smith interprets expected-utility decision models as codifying the
Humean conception of what an agent has most subjective reason to do.
His remaining argument primarily concerns what an agent has some
objective reason to do. Transferring a conclusion between those targets
would change both the evidential and comparative question. Similarly,
the statement that deliberation is rational relative to evidence does
not establish that the evidence reveals the truth
\parencite[p.~695]{Smith2021MoralReasons}.

\begin{hbexample}{def:p11-r122-selfish}{The source's entirely selfish agent}
Consider an agent whose sole intrinsic end is personal welfare. Stipulate
a situation in which acting impartially would not advance that end, and
in which the relevant facts are correctly disclosed. On the selected
Humean account, no motivation founded on that end to act impartially
survives the stipulated good deliberation. By the deliberation platitude,
the source concludes that the agent has no objective reason to do so.

The absence of a relevant desire and of a connection to its satisfaction
is essential. If impartial action brings rewards, avoids punishment,
protects a valued relationship or satisfies an additional concern, the
Humean can recognize a reason. The example does not claim that actual
people normally have only one desire or that every moral action reduces
personal welfare \parencite[p.~695]{Smith2021MoralReasons}.
\end{hbexample}

The criticism becomes a question about moral reasons only once one also
says whether the selfish agent is morally required to act impartially.
The motivational model alone does not settle that premise. Smith's
next distinction makes the required bridge explicit.

\section{Rationalism, its denial and contingent moral reasons}
\label{sec:p11-r122-rationalism}

\begin{hbdefinition}{def:p11-r122-rationalism}{Rationalism and Antirationalism in this chapter}
The source's Rationalism holds that, if moral requirements exist, their
basic requirements are knowable a priori and entail corresponding
reasons for action. Antirationalism denies that position. These labels
concern morality's knowability and its link to reasons, not merely whether
agents use reasoning.

Because Rationalism includes both commitments, its denial need not deny
each separately. A position could reject a-priori moral knowledge while
retaining a requirement-to-reason connection, or reject that connection
while taking some basic moral truths to be knowable a priori. The source's
particular Humean Antirationalist combination is more specific than the
bare negation \parencite[pp.~695--696]{Smith2021MoralReasons}.
\end{hbdefinition}

The apparent conflict between Humeanism and Rationalism has the following
premises. The selfish agent remains subject to a moral requirement despite
having no intrinsic desire served by compliance. Rationalism makes that
requirement entail a reason to comply. Humeanism, with the source's
counterfactual-deliberation account, denies any such reason under the
stipulated desire and satisfaction conditions. Those conclusions conflict.
One cannot reconcile them while retaining every premise unchanged.

This is a conditional argument, not a proof that Humeanism is false
independently of a moral theory. If one denies Rationalism, or denies that
the agent is subject to the stipulated moral requirement, that particular
contradiction is not established. Smith accordingly says the selfish
case is not automatically an objection to Humean reasons
\parencite[p.~695]{Smith2021MoralReasons}.

Smith invokes Kant for Rationalism and Brink and Railton for its denial,
and Williams and Schroeder for the Humean objective-reason interpretation.
These locate the 2021 discussion; they do not identify every argument
in those distinct works with the short chapter's reconstruction.

\begin{hbexample}{def:p11-r122-rule-cases}{The poker magician and the third-person invitation}
The source's Antirationalist analogies concern rules that apply without
necessarily supplying a reason to comply. A magician secretly engaged
to entertain poker players through card tricks need not have a reason
to shuffle as the game's rules direct. A maverick receiving an invitation
written in the third person need not have a reason to reply in the
third person under the etiquette convention being illustrated.

These are hypothetical contrasts of a rule's application and its
normative force. They do not establish that card-game violations never
involve deception or that a person can disregard every promise or social
expectation. A circumstance supplying an independent moral reason would
change the case. The etiquette rule's current prevalence is not needed
for the example's conceptual role
\parencite[pp.~695--697]{Smith2021MoralReasons}.
\end{hbexample}

Knowledge of game and etiquette rules can depend on contingent practices
and change with them. Humean Antirationalists ask why moral requirements
must have a fundamentally different kind of knowability or authority.
On their combination, moral reasons arise when complying with the correct
impartial principles serves some intrinsic desire. They can disagree
about the empirical grounds of the correct moral conception while agreeing
that an agent lacking every relevant desire has no moral reason.

Robustness does not remove that contingency. Someone intrinsically
desiring to comply with moral requirements can retain a reason to do so
even when compliance serves none of their other ends. Someone without
that intrinsic desire may have a reason only through a further connection
to another end. The first person's moral concern is more robust across
circumstances, but remains contingent on possessing that intrinsic desire.
Humeanism therefore permits genuine moral concern without making it a
rational requirement on all agents
\parencite[p.~696]{Smith2021MoralReasons}.

\section{Moral responsibility as the challenge to Humean Antirationalism}
\label{sec:p11-r122-responsibility}

Smith first grants the combination's attractions. It is economical about
psychology, using familiar belief and desire roles, and modest about norms,
assimilating moral requirements to other practical standards. His objection
is that this economy may not fit the normative conditions under which
people are appropriately held morally responsible.

The knowledge side of the objection starts from a responsibility
commitment: ordinary reasoning capacities should put basic moral
knowledge within reach regardless of a person's birthplace or historical
period. The chapter compares this with access to basic mathematical
knowledge. It expressly acknowledges that the comparison does not prove
a-priori moral knowability. It instead challenges empirical accounts tying
basic moral facts exclusively to particular human conditions. Its further
extraterrestrial example assumes that intelligent nonhuman agents could
know basic moral facts; under that assumption, those facts cannot depend
solely on distinctive facts about humans or Earth
\parencite[p.~696]{Smith2021MoralReasons}.

Neither widespread availability nor the extraterrestrial supposition alone
establishes the conclusion. An a-posteriori truth can be accessible in
many environments, and the relevant membership class remains a moral
theory's question. These are limitations on the argument, not empirical
claims about actual extraterrestrial knowledge. The issue is what an
adequate account of moral knowledge must explain.

\begin{hbdefinition}{def:p11-r122-excuses}{Exemptions, excuses and lack of concern}
An exemption in the chapter's responsibility discussion removes an agent
from the relevant ordinary competence conditions, as with infancy or a
serious absence of the reasoning capacities needed for basic moral
knowledge. An excuse can instead concern nonculpable ignorance that
one's particular circumstances trigger a moral requirement. Having no
intrinsic desire served by compliance is, on Smith's responsibility
premise, neither an exemption nor an excuse.

These distinctions concern fitting moral assessment, not a legal test,
a clinical classification, or an inference that ignorance is always
blameless \parencite[pp.~696--697]{Smith2021MoralReasons}.
\end{hbdefinition}

\begin{hbexample}{def:p11-r122-callous}{The source's callous nonhelper}
An agent is morally required to help someone in need, has the relevant
capacities and awareness, and lacks an exemption or excuse. The agent
does not help simply because the agent does not care. Smith treats this
as an appropriate candidate for blame even if no existing intrinsic
desire would be served by helping.

The requirement, the absence of excuses and the fittingness of blame
are normative premises. A report that observers blame a real nonhelper
would not by itself establish any of them. Adding a desire for approval
or fear of sanction would also change the Humean no-reason premise
\parencite[pp.~696--698]{Smith2021MoralReasons}.
\end{hbexample}

The action side of the objection can now be reconstructed in full.
First, fitting blame for failing to perform an act presupposes a reason
to perform it; Smith suggests the stronger requirement of a decisive
reason. Second, failing a moral requirement without excuse or exemption
can merit blame, including failure through callous indifference. Third,
mere absence of the relevant intrinsic desire supplies neither excuse
nor exemption. Fourth, Humeanism can make that very agent lack any reason
to comply. Retaining the first three premises makes the fourth an
inadequate account of the reason needed for fitting responsibility.

This is stronger than saying that people happen to blame nonhelpers.
It relies on what blame is appropriate and what that appropriateness
presupposes. A defender of Humeanism can contest the blame-to-reason
premise, the decisiveness claim or the applicability of the moral
requirement. Smith's argument identifies the cost of those responses;
it does not settle them by an observed majority reaction. Failures of
game rules or etiquette ordinarily lack the same implication for fitting
blame unless they also violate a moral requirement, which is the source's
reply to the earlier analogies
\parencite[pp.~696--697]{Smith2021MoralReasons}.

The proposed repair preserves the deliberation platitude but revises its
account of practical rationality. An intrinsic desire need not be true
or false to be assessable under another rational standard. Smith next
contrasts two ways of making that assessment possible.

\section{Reasons Primitivism: normative assessment of attitudes}
\label{sec:p11-r122-primitivism}

\begin{hbdefinition}{def:p11-r122-primitive}{Counting in favor as a primitive relation}
In the Reasons Primitivism Smith attributes to Parfit and Scanlon,
considerations stand in a primitive normative relation of counting in
favor of judgment-sensitive attitudes. The relation is not generally
analyzed as support for truth. Beliefs, intrinsic desires, intentions,
trust, admiration and fear belong to the relevant attitude class.

Reasons for believing do support the truth of what is believed, but the
theory treats this as a substantive feature of that subclass, not an
analysis of reasons in general. A reason to admire or intrinsically
desire something does not have to show that the admired or desired
content is true \parencite[pp.~697--698]{Smith2021MoralReasons}.
\end{hbdefinition}

The primitivist does not deny the truth relevance of evidence. The denial
is that truth-supportingness exhausts the normative reason relation.
Asked what counting in favor consists in, the account identifies it as
providing a reason, and conversely. That terminological interchange
expresses the proposed stopping point of analysis. It is not a further
argument for the relation's existence or an explanation of why a
particular consideration is a sufficient reason.

On this account, ideally rational agents acquire, retain or relinquish
beliefs and intrinsic desires in response to what they take to be
sufficient reasons for those attitudes. This introduces a parallel
between the assessment of belief and the assessment of desire, without
making the desired outcome itself an assertion of fact. For objective
reasons, the third deliberation condition adds an essential constraint:
the considerations taken to be sufficient reasons must really be
sufficient reasons. Reliable disclosure must include the normative
facts, not just ordinary causal facts
\parencite[pp.~697--698]{Smith2021MoralReasons}.

The resulting action account says that an agent has a reason to do what
would satisfy intrinsic desires the agent has sufficient reason to have,
rather than whatever would satisfy any actual intrinsic desire. In the
source's version, actions are not themselves judgment-sensitive attitudes
and so are not relata of the primitive relation. Action reasons are
derived through reasons for believing and intrinsically desiring, together
with their motivational consequences. Smith calls this a polysemy of
``reason'': a primitive attitude relation and a derived action relation
must not be conflated merely because the same noun is used
\parencite[p.~698]{Smith2021MoralReasons}.

\begin{hbdefinition}{def:p11-r122-primitive-moral}{Instrumental and noninstrumental moral reasons on the primitivist account}
An agent has a moral reason when good deliberation, with normative facts
properly disclosed, would motivate impartial action under the correct
conception of impartiality. The moral reason is noninstrumental when
the intrinsic desire there is sufficient reason to have itself has
impartial content. It is instrumental when impartial action happens to
serve a justified intrinsic desire with another content.

The source locates robust moral reasons in the noninstrumental case:
they do not depend on some differently directed justified desire
happening to be advanced by impartial action
\parencite[p.~698]{Smith2021MoralReasons}.
\end{hbdefinition}

Smith proposes combining this with Rationalism by understanding moral
requirements through decisive noninstrumental moral reasons. For a
specified required act, there must be sufficient reason for a suitably
impartial intrinsic desire; that act must satisfy the desire; and the
reason must be strong enough, relative to other reasons, that good
deliberation would produce the strongest motivation to perform it.
The combination thereby supplies the reason which the responsibility
argument says the callous agent possesses. Actual lack of concern does
not remove the reason to have the appropriate concern
\parencite[pp.~698--699]{Smith2021MoralReasons}.

This is a possible integration of the two positions, not an inference
from primitive normativity alone to every moral requirement. One still
needs the correct impartial content and the sufficiency and decisiveness
premises. The source's earlier a-priori-knowledge commitment is also part
of Rationalism; calling the relation primitive does not, by itself,
establish how agents know the relevant basic normative facts.

The source gives two linked objections. A primitive relation may be an
unnecessarily costly metaphysical posit. Its psychological demands are
also substantial: attitudes must track judgments about that relation,
and a person denying that such a relation exists appears irrational by
the account's own ideal. This is Smith's criticism of the theory's
commitments, not a measured finding that primitivists are cognitively
impaired or that actual people cannot have norm-sensitive attitudes.
If a less demanding account preserves the intended responsibility
connection, the source argues that parsimony and plausibility favor
considering it \parencite[p.~699]{Smith2021MoralReasons}.

\section{Constitutivism: protecting the capacities for deliberation}
\label{sec:p11-r122-constitutivism}

Smith's anti-Humean Constitutivism is the less radical alternative he
associates with Korsgaard and his own distinct 2013 treatment. It keeps
the Humean principles where they are adequate and adds principles
governing intrinsic desires. Belief and desire retain their different
truth roles. The additional criticism of desire concerns what is
constitutive of deliberating well, not whether the desired state is true.

The first step is a functional-role premise: the full rational principles
for belief and desire govern the roles these states play in good
deliberation. Those roles are supposed to make the states the psychological
kinds they are, rather than being external advice added after describing
them. If having particular intrinsic desires is necessary for
that deliberation, the account makes those desires rationally required.
It therefore denies the Humean inference from a desire's nonassertoric
content to the absence of all rational standards for it
\parencite[p.~699]{Smith2021MoralReasons}.

\begin{hbdefinition}{def:p11-r122-help-noninterference}{The source's desires to help and not interfere}
The proposed additional intrinsic desires are to avoid interfering
with anyone's exercise of deliberative capacities, oneself included,
and to do what one can to ensure that everyone has such capacities
to exercise. They concern possession and exercise of capacities, not
just success at one isolated choice. Their contents are impartial in
including one's own and other agents' capacities under the same concern
\parencite[p.~699]{Smith2021MoralReasons}.
\end{hbdefinition}

The argument for these desires begins with robustness. Good deliberation
cannot depend merely on a lucky episode of disclosure; the relevant
world must be reliably available. Deliberation is undermined by others'
deception, self-deception, coercion by others or oneself, and preventable
deterioration of capacities that oneself or others leave unaddressed.
Reliable dispositions concerning these conditions matter as well as
their accidental absence on a particular occasion. Smith moves from
these conditions to a social context of good deliberators whose strong
intrinsic desires sustain everyone's capacities and protect their exercise
\parencite[p.~699]{Smith2021MoralReasons}.

The source's inference has a substantive dependency. Conditions favorable
to reliable deliberation do not simply constitute proof that each agent
must intrinsically desire their protection, still less universal impartial
protection. A person might secure reliable information for themselves
while obstructing someone else's deliberation. That possibility alone
does not refute the constitutivist ideal, but shows why a merely
instrumental security argument is insufficient. The account must defend
the stronger claim that the specified intrinsic concern, including its
social and impartial scope, is constitutive of good deliberation.
Treating it as a constitutive condition makes the later inference
conditional on that claim; it does not establish the claim by renaming it.

Given that constitutive premise and the deliberation platitude, agents
have noninstrumental moral reasons to help and not interfere independently
of their other actual ends. The source identifies the corresponding
moral requirements through those reasons and says the account entails
Rationalism and fits moral responsibility. Its callous person can
therefore possess a reason to help while lacking the actual desire that
good deliberation would require. This reconstruction retains the
premises needed for that result: the constitutive status of the desires,
their impartiality, the ideal-motivation account of objective reasons,
and the proposed moral requirement/reason connection
\parencite[pp.~699--700]{Smith2021MoralReasons}.

The claim to the full Rationalist position additionally owes its
a-priori-knowability commitment. The short chapter does not give a
separate demonstration that ordinary agents know the basic constitutive
or moral facts a priori. Nor does its robustness discussion determine
a complete ranking of every conflict between capacity protection and
another end. Those are limitations of the argument supplied, not
reasons to suppress its conditional conclusion.

\begin{hbconvention}{def:p11-r122-constitution-limits}{Permissible other desires and the selected moral conception}
Constitutivism in this source leaves room for pursuing other intrinsic
desires so long as doing so is compatible with acting on the moral
reasons to help and not interfere when those reasons are decisive.
It is therefore not a proposal to replace every personal project with
constant assistance to everyone.

The chapter associates this with a deontological interpretation protecting
each person's freedom to choose a life consistently with others' equal
freedom. On that particular interpretation, welfare is not intrinsically
morally significant: it matters morally insofar as it affects freedom.
This is the source's substantive proposal, not a consequence shared by
all impartial moral theories or all forms of Constitutivism
\parencite[p.~700]{Smith2021MoralReasons}.
\end{hbconvention}

Calling the alternative less radical refers to its retention of ordinary
belief/desire roles and reasons to pursue compatible personal ends. It
does not establish that its added constitutive premise is undisputed,
that it has better empirical fit, or that the primitive-relation rival
has thereby been conclusively refuted.

\section{Comparison with reasons, structure and economic choice}
\label{sec:p11-r122-comparisons}

Broome's selected treatment at \cref{sec:p11-reasons-and-rationality}
distinguishes motivating reasons, normative reasons, coherent attitudes
and successful performance. Smith's account does not collapse those
items: a person's actual concern can motivate conduct, while the
relevant normative reason is analyzed through a counterfactual concerning
good deliberation. A callous agent's failure to be motivated therefore
does not settle the normative question on the anti-Humean views.

The sources ask different reduction questions. Broome disputes identifying
the person's mental and structural rationality with correct response to
actual reasons; his insurance case separates mental equality from
worldly success. Smith asks what reasons an agent has, using ideal
motivation under reliably disclosed facts, and disputes particular
accounts of the rational evaluation of intrinsic desire. His objective
idealization should not be inserted as Broome's definition of the actual
person's degree of rationality. Broome's distinction between a standard
and a reason to conform also illuminates Smith's poker/etiquette contrast,
without proving Smith's premise that fitting moral blame requires a
decisive reason \parencite[pp.~693--697]{Smith2021MoralReasons}.

TSR at \cref{sec:p11-rationality-10-5} organizes acts through intended
patterns, continuing persons and interpersonal practices of reasons.
Smith's Constitutivism also has a social and noninstrumental ambition,
but grounds its additional desires in the conditions of good deliberation.
These are not the same constitutive claim. A promise's place in an
accepted continuing practice does not automatically establish that
everyone's intrinsic desire to preserve deliberative capacities is
necessary for rational agency. Conversely, a capacity-protection reason
does not by itself validate every established role or promised act.

The economic comparison at \cref{sec:p19-rationality-9-4} keeps
preference order and expected-utility assumptions distinct from
self-regarding motivation, emotions, computation and empirical fit.
Smith's critique of treating intrinsic desires as exempt from rational
assessment concerns a philosophical interpretation of what the choice
model leaves fixed. Neither an ordinal representation nor maximization
of a supplied expected utility establishes an end's moral adequacy.
Introducing an other-regarding utility can model a concern; it does
not prove that an agent has a noncontingent reason to possess it.

These distinctions also constrain comparisons with deontic reasoning.
A derivation from supplied obligation premises can test a consequence
relation; it does not establish the premises' moral authority or explain
which desires should guide a reasoner. Smith's moral and agency premises
cannot be supplied merely by adding an obligation symbol to a description
of observed behavior.

\section{A bounded use in psychological inquiry}
\label{sec:p11-r122-use}

\begin{hbexample}{def:p11-r122-comparison-case}{One helping act and three explanations}
Consider a stipulated agent who can give needed help. Suppose helping
is morally required, and separately record the agent's factual beliefs,
intrinsic concerns, instrumental beliefs, actual motivation and action.
A desire for approval can make helping an instrumentally motivated
option. An intrinsic desire to comply with impartial requirements can
make its motivation robust across changes in approval. Neither fact
alone establishes that the desire is rationally required.

The Humean explanation asks whether helping would satisfy an actual
intrinsic desire under correct factual disclosure. A primitivist
explanation asks whether good deliberation would supply a sufficiently
justified intrinsic concern with the relevant content. The source's
constitutivist explanation asks whether helping expresses the intrinsic
concerns required to sustain and protect deliberative capacities.
The same observed helping act can fit more than one of these accounts.

Now stipulate a callous nonhelper with no relevant intrinsic desire,
normal capacities, relevant knowledge and no excuse. Under Smith's
responsibility premises, fitting blame requires a reason which the
selected Humean account cannot supply. The anti-Humean accounts can
attribute that reason through an ideally required concern. This is the
source's philosophical contrast applied to a declared case, not a
reported population result or a measurement of ideal motivation.
\end{hbexample}

\begin{hbmethod}{def:p11-r122-use}{Keep the moral and psychological assumptions inspectable}
Identify the proposed moral conception, its scope and interpretation of
impartiality. State separately whether basic moral facts are taken to
be knowable a priori and whether moral requirements entail reasons.
Do not use the single label Antirationalism to conceal which commitment
is rejected.

Declare the theory of practical rationality. Is the agent's actual
intrinsic-desire set taken as given, evaluated through sufficient
primitive attitude reasons, or constrained by constitutive requirements
of deliberation? Specify whether the conclusion concerns some or most
reason, subjective or objective reason, actual motivation or ideal
counterfactual motivation. Record what reliable disclosure is supposed
to reveal, including normative facts if the theory requires them.

For responsibility, distinguish an observed blaming response from a
judgment that blame is fitting. Specify exemptions, excuses and the
premise linking fitting blame to reasons. For an empirical comparison,
specify the task, recruitment, incentives, measures, design, uncertainty
and alternative explanations under \cref{def:p19-evidence-record}.
Do not pretend that an intrinsic desire, an ideal reason or a normative
requirement has been measured merely because a participant chose an act
or endorsed a questionnaire statement.

A successful application separates the callous-agent premises and shows
exactly why the accounts give different verdicts. A failed application
infers noncontingent moral reasons from observed cooperation, treats a
desire's lack of truth conditions as proof that no rational criticism
is possible, or declares universal impartial desire necessary merely
because supportive social conditions improve deliberation.
\end{hbmethod}

The chapter's conclusion recommends that psychologists and collaborating
philosophers make these commitments explicit, then investigate whether
some combinations yield more productive research programs. That is a
proposal for inquiry, not a reported comparative outcome. Better
prediction, if established in a specified task, would support a
descriptive model there; it would not by itself prove a-priori moral
knowledge, a primitive normative ontology, or the authority of all
constitutive norms \parencite[p.~700]{Smith2021MoralReasons}.

\begin{proofboundary}{def:p11-r122-boundary}{The standing of the argument}
The supplied argument includes the source's moral characterization,
Humean psychological account, ideal-deliberation distinctions, contingent
reason examples, responsibility challenge, both anti-Humean alternatives,
their costs and the final research recommendation. Its normative
premises and constitutive bridges remain explicit rather than being
promoted to formal or empirical results.

There is no source formal theorem or rule requiring a proof annex and
no empirical study whose sample or effect could be reconstructed from
this chapter. Its references to other works identify philosophical
provenance and fuller treatments, not a claim that those whole works
have been supplied. Comparisons with Broome, TSR and economic
representation preserve their respective questions and assumptions.
The argument permits disagreement about its premises; it does not
manufacture a settled answer by treating a proposed moral norm as a
definition of universally accepted rationality.
\end{proofboundary}

\index{moral reasons!Smith}
\index{impartiality!moral principles}
\index{Rationalism!moral requirements}
\index{Antirationalism!moral requirements}
\index{Humeanism!intrinsic desires}
\index{Reasons Primitivism}
\index{Constitutivism!deliberative capacities}
\index{moral responsibility!reasons and blame}
